%\documentclass[12pt,oneside]{amsart}
\documentclass{scrartcl}

%\documentclass{article}

\usepackage{amsthm,amsmath,amssymb}
\usepackage[numbers]{natbib}

%\usepackage[colorlinks]{hyperref}
\usepackage{hyperref}
\hypersetup{
    colorlinks,%
    citecolor=black,%
    filecolor=black,%
    linkcolor=black,%
    urlcolor=black
}
\usepackage{hypernat}
\usepackage[top=2.5cm, bottom=2.5cm, left=2.8cm,
right=2.8cm]{geometry}
\usepackage{graphicx}


%\setlength{\parskip}{0pt}
%\setlength{\parsep}{0pt}
%\setlength{\headsep}{0pt}
%\setlength{\topskip}{0pt}
%\setlength{\topmargin}{0pt}
%\setlength{\topsep}{0pt}
%\setlength{\partopsep}{0pt}

\linespread{1}

%\usepackage{marvosym}

%\textwidth = 6.5 in \textheight = 9.2 in \oddsidemargin = 0.0 in
%\evensidemargin = 0.0 in \topmargin = -0.0 in \headheight = 0.0 in
%\headsep = 0.0in \parskip = 0.0in \parindent = 0.0in
%\def\baselinestretch{0.871}


\newtheorem{theorem}{Theorem}[section]
\newtheorem{proposition}[theorem]{Proposition}
\newtheorem{problem}[theorem]{Problem}
\newtheorem{fact}[theorem]{Fact}
\newtheorem{lemma}[theorem]{Lemma}
\newtheorem{defn}[theorem]{Definition}
\newtheorem{corollary}[theorem]{Corollary}
\newtheorem{remark}{Remark}[section]
\newtheorem{example}[theorem]{Example}

\newcommand{\E}{{\mathbb E}}
\newcommand{\se}{{\mathcal{E}}}
\newcommand{\R}{{\mathbb R}}
\renewcommand{\P}{{\mathbb P}}
\newcommand{\ac}{{\mathcal{A}}}
\newcommand{\A}{{\mathcal{A}}}
\newcommand{\B}{{\mathcal{B}}}
\newcommand{\C}{{\mathcal{C}}}
\newcommand{\M}{{\mathcal{M}}}
\newcommand{\Mf}{{\mathfrak{M}}}
\newcommand{\T}{{\mathcal{T}}}
\newcommand{\Z}{{\mathcal{Z}}}
\newcommand{\K}{{\mathcal{K}}}
\newcommand{\lc}{{\mathcal{L}}}
\newcommand{\Ps}{{\mathcal{P}}}
\newcommand{\G}{{\mathcal{G}}}
\newcommand{\pa}{{\mathring{p}}}
\newcommand{\F}{{\cal F}}
\newcommand{\samp}{{\mathcal{X}}}
\newcommand{\X}{{\mathcal{X}}}
\newcommand{\hi}{{\hat{\Phi}}}
\newcommand{\data}{{\text{data}}}
\newcommand{\relu}{{\text{ReLU}}}

\newcommand{\ridge}{{\hat{\beta}_{\text{ridge}}(\lambda)}}
\newcommand{\lasso}{{\hat{\beta}_{\text{lasso}}(\lambda)}}
\newcommand{\ridgemu}{{\hat{\mu}_t^{\text{ridge}}(\lambda)}}
\newcommand{\lassomu}{{\hat{\mu}_t^{\text{lasso}}(\lambda)}}


\newcounter{rcnt}[section]
\renewcommand{\thercnt}{(\roman{rcnt})}

\def\qt#1{\qquad\text{#1}}

\def\argmin{\mathop{\rm argmin}}
\def\argmax{\mathop{\rm argmax}}


\setlength{\parskip}{1.4 \medskipamount}
%\sloppy
%\linespread{1.3}

\begin{document}

\title{STAT 238 - Bayesian Statistics  \\
  Lecture Twenty Nine} 
\subtitle{Spring 2026, UC Berkeley}

\author{Aditya Guntuboyina}


\date{10 April 2026}

\maketitle

\section{Gibbs Sampler for Mixture Models}

\subsection{Known variances and proportion}

We observe real-valued data $y_1, \dots, y_n$ and consider the model:
\begin{align}\label{mixmod}
  y_i \overset{\text{i.i.d}}{\sim} (1 - w) N(\mu_0, 1) + w N(\mu_1, 1)
\end{align}
with $w$ known and fixed (e.g., $w = 0.3$ or $0.7$ as in the simulations) and
$\mu_0, \mu_1$ being unknown. We will describe the Gibbs sampler
algorithm in this setting first, and then generalize it to the case
where $w$ is unknown and we have variance parameters $\sigma_0^2$ and
$\sigma_1^2$. 

The log-likelihood corresponding to \eqref{mixmod} is:
\begin{align*}
  \sum_{i=1}^n \log \left(p \phi(y_i, \mu_1, 1) + (1 - p) \phi(y_i,
  \mu_2, 1) \right), 
\end{align*}
where $\phi(x, \mu, \sigma^2)$ is the normal density with mean $\mu$,
variance $\sigma^2$ evaluated at $x$.


We will take a standard prior for $\mu_0, \mu_1$ e.g., $\mu_0, \mu_1
\overset{\text{i.i.d}}{\sim} N(0, C)$ for a large $C$. The posterior
of $\mu_0, \mu_1$ is then:
\begin{align*}
\pi(\mu_0, \mu_1 \mid \text{data}) \propto \phi(\mu_0, 0, C) \phi(\mu_1,
  0, C) \prod_{i=1}^n \left((1 - w) \phi(y_i, \mu_0, 1) + w \phi(y_i,
  \mu_1, 1) \right). 
\end{align*}
This posterior cannot be evaluated in closed form and numerical
methods need to be used. A standard approach is to use the Gibbs
sampler with augmentation. First observe that the model \eqref{mixmod}
can be rewritten in the following way: 
\begin{align*}
  z_i \overset{\text{i.i.d}}{\sim} \text{Bernoulli}(w) ~~ \text{ and }
  ~~ y_i \mid z_i = 1 \sim N(\mu_1, 1) ~~ \text{ and } ~~ y_i \mid z_i
  = 0 \sim N(\mu_0, 1). 
\end{align*}
It should be clear that, under the above model, the marginal
distribution of $y_i$ coincides with \eqref{mixmod}. $z_1, \dots, z_n$
can be thought of as unobserved latent variables which represent which
of the two populations (corresponding to the distributions $N(\mu_0,
1)$ and $N(\mu_1, 1)$ respectively) the observation $y_i$ comes from.

Gibbs sampler is implemented for jointly sampling from the posterior
of $\mu_0, \mu_1, z_1, \dots, z_n$ given the data. This requires being
able to sample from the full conditionals
\begin{align*}
  z \mid \mu_0, \mu_1, y ~~ \text{ and } ~~ (\mu_0, \mu_1) \mid z, y. 
\end{align*}
where $y = (y_1, \dots, y_n)$ is the data and $z = (z_1, \dots,
z_n)$. It is easy to see that these full conditionals can be written
in closed form as follows. Given $\mu_0, \mu_1, y_1, \dots, y_n$, the
variables $z_1, \dots, z_n$ are independent with
\begin{align*}
  z_i \mid \mu_0, \mu_1, y \sim \text{Bernoulli} \left(\frac{w
  \phi(y_i, \mu_1, 1)}{w \phi(y_i, \mu_1, 1) + (1 - w) \phi(y_i,
  \mu_0, 1)}\right). 
\end{align*}
On the other hand, given $z_1, \dots, z_n, y_1, \dots, y_n$, the
variables $\mu_0, \mu_1$  are independent with
\begin{align*}
  \mu_1 \mid z, y \sim N \left(
  \frac{\sum_{i=1}^n y_i z_i}{n_1 + (1/C)}, \frac{1}{n_1 + (1/C)}
  \right) \text{ and } \mu_0 \mid z, y \sim N \left(
  \frac{\sum_{i=1}^n y_i (1-z_i)}{n_0 + (1/C)}, \frac{1}{n_0 + (1/C)}
  \right) 
\end{align*}
where $n_0$ is the number of $z_i$'s that are equal to 0, and $n_1$ is
the number of $z_i$'s that are equal to 1. Note that in the limit as
$C \rightarrow \infty$, we can write 
\begin{align*}
  \mu_1 \mid z, y \sim N \left(\frac{\sum_{i=1}^n y_i
  z_i}{\sum_{i=1}^n z_i}, \frac{1}{\sum_{i=1}^n z_i} \right) ~~ \text{
  and } ~~   \mu_0 \mid z, y \sim N \left(\frac{\sum_{i=1}^n y_i
  (1 - z_i)}{\sum_{i=1}^n (1 -z_i)}, \frac{1}{\sum_{i=1}^n (1 - z_i)} \right)
\end{align*}
Based on these full conditional distributions, the Gibbs sampler
algorithm takes the following form:
\begin{enumerate}
\item Initialize $\mu_0^{(0)}, \mu_1^{(0)}$.
\item Repeat the following for $t = 0, 1, 2, \dots$:
  \begin{enumerate}
  \item Generate $z_1^{(t)}, \dots, z_n^{(t)}$ via:
    \begin{align*}
      z_i^{(t)} \sim \text{Bernoulli} \left(\frac{w
  \phi(y_i, \mu^{(t)}_1, 1)}{w \phi(y_i, \mu^{(t)}_1, 1) + (1 - w) \phi(y_i,
  \mu^{(t)}_0, 1)}\right). 
    \end{align*}
  \item Generate $\mu_0^{(t+1)}, \mu_1^{(t+1)}$ via:
\begin{align*}
&  \mu^{(t+1)}_1  \sim N \left(\frac{\sum_{i=1}^n y_i
  z^{(t)}_i}{\sum_{i=1}^n z^{(t)}_i}, \frac{1}{\sum_{i=1}^n z^{(t)}_i}
  \right) \\ & \mu^{(t+1)}_0  \sim N \left(\frac{\sum_{i=1}^n y_i
  (1 - z^{(t)}_i)}{\sum_{i=1}^n (1 -z^{(t)}_i)}, \frac{1}{\sum_{i=1}^n
  (1 - z^{(t)}_i)} \right).
\end{align*}    
  \end{enumerate}
\end{enumerate}
As we saw in the simulations in the last couple of lectures,
initialization is very important. The log-likelihood can have multiple
modes only one of which is the correct mode (in the sense of having
large likelihood). If the Gibbs sampler is not properly initialized,
the algorithm would sample from a spurious peak.

\subsection{Unknown variances and proportion}

Now consider the model:  
\begin{align*}
  y_i \overset{\text{i.i.d}}{\sim} (1 - w) N(\mu_0, \sigma_0^2) + w
  N(\mu_1, \sigma_1^2). 
\end{align*}
The unknown parameters are now $w, \mu_0, \mu_1, \sigma^2_0,
\sigma_1^2$. We place the following independent priors on these
variables: 
\begin{align}\label{priors}
  w \sim \text{Beta}(a, b)~~ \text{ and } ~~ \mu_0, \mu_1 \sim N(m,
  s^2) ~~ \text{ and } ~~
  \sigma_0^2, \sigma_1^2 \sim IG(\alpha, \beta). 
\end{align}
The proportion parameter $w$ is constrained to lie in $[0, 1]$ so it
is natural to take the Beta prior for it. A standard choice here is $a
= b = 1$ (which corresponds to the uniform prior). For $\mu_0, \mu_1$,
we are using the $N(m, s^2)$ prior. Standard choice is $m = 0$ and $s$
very large. The Inverse Gamma prior $IG(\alpha, \beta)$ corresponds to
the density:
\begin{align*}
  \propto x^{-\alpha - 1} \exp(-\beta/x) I\{x > 0\}. 
\end{align*}
The standard uninformative prior for $\sigma$ is $\propto \sigma^{-1}
I\{\sigma > 0\}$. The corresponding prior density for $\sigma^2$
is also $\propto x^{-1} I\{x > 0\}$. This is a special case of
$IG(\alpha, \beta)$ corresponding to $\alpha = \beta = 0$.

The prior \eqref{priors} is therefore uses conjugate families while
including the standard uninformative choices as special cases. The
reason for conjugacy is that the full conditional distributions
corresponding to the posterior distribution can be written in closed
form. We shall look at the formulae in the next lecture. 

%\begin{thebibliography}{99}

%\bibitem[Roberts and Rosenthal(2004)]{RobertsRosenthal2004}
%Roberts, G. O. and Rosenthal, J. S. (2004).
%General state space Markov chains and MCMC algorithms.
%\textit{Probability Surveys}, \textbf{1}, 20--71.

%\bibitem[Meyn and Tweedie(2009)]{MeynTweedie2009}
%Meyn, S. P. and Tweedie, R. L. (2009).
%\textit{Markov Chains and Stochastic Stability}, 2nd ed.
%Cambridge University Press.


  
%\bibitem[Chib and Greenberg(1995)]{chib1995}
%Chib, S. and Greenberg, E. (1995).
%Understanding the Metropolis--Hastings algorithm.
%\textit{The American Statistician}, \textbf{49}, 327--335.


%\bibitem[MacKay and Peto(1995)]{MacKay1995HierarchicalDirichlet}
%D.~J.~C. MacKay and L.~C. Peto,
%\newblock ``A hierarchical Dirichlet language model,''
%\newblock \emph{Natural Language Engineering}, vol.~1,
%no.~3, pp.~289--308, 1995.

%\bibitem[Rubin(1981)]{Rubin1981BayesianBootstrap}
%D.~B. Rubin,
%\newblock ``The Bayesian bootstrap,''
%\newblock \emph{The Annals of Statistics}, vol.~9,
%no.~1, pp.~130--134, 1981.
  
%\bibitem[Fisher(1934)]{Fisher1934}
%R.~A. Fisher,
%\newblock ``Probability, likelihood and quantity of information in the logic of
%uncertain inference,''
%\newblock \emph{Proceedings of the Royal Society of London. Series~A,
%Containing Papers of a Mathematical and Physical Character}, vol.~146,
%no.~856, pp.~1--8, 1934.

%\bibitem[Efron(2010)]{Efron}
%Efron, B. (2010)
%\textit{Large-scale inference: empirical Bayes methods for estimation,
%testing, and prediction}. Cambridge University Press, 2010.

%\bibitem[Shumway and Stoffer(2017)]{ShumwayStoffer}
%Shumway, R. H. and Stoffer, D. S. (2017)
%\textit{Time Series Analysis and Its Applications: With R
%  Examples}. Springer, 4th edition. 
  
%     \bibitem[Jaynes(2003)]{Jaynes} Jaynes, E. T, (2003)
%  \textit{Probability theory: the logic of science}. Cambridge
%  University Press, 2003.
 
%  \bibitem[Sinai(1992)]{Sinai} Sinai, Y. G, (1992)
%  \textit{Probability theory: an introductory course}. Springer
%  Verlag, 1992.  

%\bibitem[{deGroot(1986)}]{DeGrootBlackwell}DeGroot, M. H. (1986)
%       A conversation with David Blackwell.
%\newblock \emph{Statistical Science}, \textbf{1}, 40--53.

%         \bibitem[Mosteller(1987)]{Mosteller} Mosteller, F, (1987)
%  \textit{Fifty challenging problems in probability with
%    solutions}. Courier Corporation, 1987.

%    \bibitem[MacKay(2003)]{MacKay} MacKay, D. J. C., (2003)
%  \textit{Information theory, inference, and learning
%  algorithms}. Cambridge University Press, 2003.

%\bibitem[{Lindley(2013)}]{Lindley}Lindley, D. V. (2013)
%   \textit{Understanding Uncertainty}. John Wiley and sons, 2013.


%\end{thebibliography}







\end{document}

%%% Local Variables:
%%% mode: latex
%%% TeX-master: t
%%% End:
